\chapter{NGUỒN DỮ LIỆU VÀ RULE ẢNH HƯỞNG ĐẾN CÁCH HIỂU SỐ LIỆU}

\section{Nguồn dữ liệu chính theo góc nhìn người đọc report}

Ở mức reader-facing, các nguồn cần nhớ trước tiên là:

\begin{enumerate}[label=-]
\item \texttt{energy\_kpi}: nguồn chính thức cho KPI reporting và official energy totals;
\item \texttt{utility\_usage}: nền cho utility usage surfaces;
\item \texttt{processvalue}: nền cho sensor monitoring và các thống kê sensor-level.
\end{enumerate}

\section{Những rule ảnh hưởng trực tiếp đến cách đọc report}

Một số rule quan trọng sẽ cần được viết kỹ hơn ở các checkpoint sau:

\begin{enumerate}[label=-]
\item official totals không nên bị hiểu như số recompute tùy ý từ raw views;
\item zero là dữ liệu hợp lệ, còn missing value phải được hiểu khác zero;
\item current vs previous comparison chỉ có ý nghĩa đầy đủ khi kỳ đối chiếu phù hợp;
\item KPI dùng coverage-first logic, nên coverage status ảnh hưởng trực tiếp đến cách diễn giải kết quả;
\item weekly và monthly surfaces không đơn giản chỉ là daily phóng to, mà có các rule chọn kỳ và layout riêng.
\end{enumerate}

\section{Mục tiêu của chapter này ở bản hoàn chỉnh}

Ở các checkpoint nội dung sau, chapter này sẽ nối từng phần trong report với:

\begin{enumerate}[label=-]
\item nguồn dữ liệu tương ứng;
\item rule tổng hợp hoặc rule chọn số liệu;
\item caveat diễn giải mà người đọc báo cáo cần biết.
\end{enumerate}
